\section{Remarks}\label{sec:remarks}

\begin{rem}
Euler's original argument multiplies divergent series as if they
converged; the inequalities of \cref{lem:log} are what make it rigorous,
and \thmref{divergence} is what it was after.
\end{rem}

\begin{proof}[Proof of the claims in \cref{sec:product}]
Each claim there follows from absolute convergence.
\end{proof}

\begin{lem}\label{lem:tail}
For $s \ge 2$, $\sum_{n > x} n^{-s} \le 1/x$.
\end{lem}
The bound is crude.
\begin{proof}
Compare with $\int_x^\infty t^{-2}\,dt$.

The integral is $1/x$.
\end{proof}

\ifdefined\eulerconst We write $\eulerconst$ for Euler's constant.\fi

\subsubsection*{Further reading}
\Cref{sec:sharper} and the section \nameref{sec:remarks} point to
\citet{primer2001}; \vref{lem:tail} is used nowhere else.

\section*{Acknowledgements}
We thank the reader for reading.
